\section{The two infinity disks and their rational double zeros}
\label{sec-fixed-infinity-certificate}

We give the full ordinary certificate IF1--IF3, including the central
height identity. Put $q=1/z$, $v=Wq^3$, $v(0)=1$. Then
\[
\begin{split}
 v^2&=1-27q^2+99q^4-9q^6,\\
 R_1&=((q^{-2}-9)/4,v/(8q^3)),\qquad
 R_2=((33-9q^2)/4,-9v/8).
\end{split}
\]
Set $q=5s$ and $\mathcal A=\mathbb Z_5\langle s\rangle$, with its
complete Gauss valuation.

\paragraph{IF1: all denominators and tails on the entire disk.}
The canceled expression
$\Xi=q^{-81}\delta_1(R_1)/\delta_2(R_2)$ is an analytic unit, with
\[
 \Xi(q)=\Xi(0)(1+q^2 B(q^2)),\quad B\in\mathbb Z_5[[q^2]],\qquad
 T_i\in5\mathcal A.
\]
Here is a formal-series proof. Hensel recursion gives
$v\in1+q^2\mathbb Z_5[[q^2]]$ and
$t_1=-2q(1-9q^2)/v$. On $E_1$, the zero-slope recursion gives
$w=t^3U$, $U\in1+t^4\mathbb Z_5[[t^2]]$,
$x=t^{-2}U^{-1}$ and $y=-t^{-3}U^{-1}$.
The integral division polynomials have leading terms
$\psi_9=9x^{40}$, $\phi_9=x^{81}$ and $D=9x^{120}$;
see \cite[Lemma 5.22]{SutherlandDivisionPolynomials2023} for the first two,
and differentiate to obtain the last. Therefore
\[
 t^{80}\psi_9\in9+t^2\mathbb Z_5[[t^2]],\quad
 t^{162}\phi_9\in1+t^2\mathbb Z_5[[t^2]],\quad
 t^{240}D\in9+t^2\mathbb Z_5[[t^2]].
\]
RD1, using the integral inverse of the last unit, yields
\[
 \delta_1(t)=t^{81}V(t^2),\quad V\in1+t^2\mathbb Z_5[[t^2]],\qquad
 t([9]Q)\in t(9+t^2\mathbb Z_5[[t^2]]).
\]
Thus $q^{-81}\delta_1(R_1)$ is $(-2)^{81}$ times an even integral unit.
The second quotient is an integral even point with center $(33/4,-9/8)$.
Writing $e=q^2$, its actual first jets modulo five, and the full
division recurrences applied to them, give
\[
 (x,y)=(2+4e,2+3e),\quad
 \psi_9=3e,\quad \phi_9=4+3e,\quad \omega_9=2+e\pmod{(5,e^2)}.
\]
Consequently $\delta_2=-\phi_9/\omega_9$ is an even integral unit,
and $T_2=-\phi_9\psi_9/\omega_9$ has constant in $5\mathbb Z_5$
and remaining terms in $q^2\mathbb Z_5[[q^2]]$. Substitution $q=5s$
gives restricted series and proves the assertions. Moreover
\[
 \log_5\Xi(q)-\log_5\Xi(0)\in25\mathcal A,
\]
because its $j$th logarithm term has valuation at least
$2j-v_5(j)$, tending to infinity. This controls the entire tail.

\paragraph{IF2: the central value is exactly zero.}
The doubling law gives $2P_2=(33/4,9/8)$ for $P_2=(6,9)$, hence
$R_2(0)=-2P_2$. As $q\to0$, the singular terms combine as
\[
 \lambda_{01}(R_1)-2\log_5z
 =-2\log_5(t_1/q)-2R_{01}(t_1)\longrightarrow-2\log_5 2.
\]
Also $\log_{E_1}(R_1)\to0$. Evenness of local heights and the
rank-one quadratic identity give
\[
 F_0(0)=-2\log_5 2-\lambda_{02}(2P_2)+H_{02}(2P_2).
\]
The height difference at this actual rational point can be computed
at every place other than five. At two, $v_2(x)=-2$ and
$3x^2-189$ has negative valuation, so Cremona case (a) in the
previously fixed twice-Silverman normalization gives $2\log_5 2$.
At three the exact tests are
\[
 v_3(x)=1,\quad v_3(y)=2,\quad A=5,\quad B=2,\quad C=6,\quad v_3(c_4)=4.
\]
Thus case (c), $C\ge3B$, gives $-(4/3)\log_5 3$. At all other
primes different from five, the model has good reduction and the
coordinates are integral, so the local height is zero. Changing
the splitting affects only five. Therefore
\[
 H_{02}(2P_2)-\lambda_{02}(2P_2)
 =2\log_5 2-\frac43\log_5 3,
\]
which proves $F_0(0)=-(4/3)\log_5 3$ and $f(0)=0$.
No density of global rational points near infinity is assumed here.

\paragraph{IF3: an exact even factor.}
There is $U_\infty\in1+5\mathcal A$ such that
\[
 f(5s)=5s^2U_\infty(s).
\]
Indeed put $\ell_i=\log_{E_i}(R_i)$. The actual differential pullbacks give
\[
 \ell_1'=-2/v,\quad \ell_2'=2q/v,\quad
 \ell_1(0)=0,\quad \ell_2(0)=-2\log_{E_2}(P_2).
\]
Integration of the integral even series $1/v$ and substitution $q=5s$
give $\ell_1(5s)/5\equiv-2s\pmod{5\mathcal A}$ and
$\ell_2(5s)-\ell_2(0)\in25\mathcal A$. At every degree $j$ the
integration denominator is controlled by $j-v_5(j)\to\infty$;
the omitted first degrees are three and two respectively.
ZS puts $R_{0i}(T_i)$ and their center differences in $25\mathcal A$.
The logarithm change has this same bound by IF1. The second squared
logarithm changes by $125\mathcal A$, since its center is in
$5\mathbb Z_5$ and its change is in $25\mathcal A$; multiplication
by $\alpha_2^0$ of valuation $-1$ leaves $25\mathcal A$.
Subtract the exact center value. The only surviving term in
$g=f(5s)/5$ modulo five is
\[
 g\equiv-(5\alpha_1^0)(\ell_1(5s)/5)^2
 \equiv-4s^2=s^2\pmod{5\mathcal A},
\]
using the exact digit $5\alpha_1^0\equiv1$. DS proves exact evenness
and IF2 proves zero constant term. Dividing the restricted series by
$s^2$ therefore gives $U_\infty\in1+5\mathcal A$. Its value is a
unit at every $s\in\mathbb Z_5$. The only zero is the central rational
infinity point and its multiplicity is exactly two. DS transports
this statement to the other infinity disk.

The exact replay records the division jets, doubling coordinates and
height valuation tests. The all-tail estimates, height identity and
zero factorization above are ordinary proofs, not consequences of
finite numerical root sampling or of a completed Lean height theory.
